\documentclass[11pt]{article}
\usepackage{amsmath,amssymb,amsthm,hyperref}
\newtheorem{proposition}{Proposition}
\newtheorem{corollary}{Corollary}
\newcommand{\RR}{\mathbb R}
\newcommand{\FF}{\mathbb F}
\newcommand{\rank}{\operatorname{rank}}
\title{Weighted support bounds for separately compiled line selectors}
\author{Joint-cell label audit}
\date{November 2026}
\begin{document}
\maketitle

\paragraph{Status and scope.}
This is a written proof using an external rank theorem, not a Lean theorem.
It concerns exact rational line selectors compiled separately before decoding.
It does not bound arbitrary xorformer scores, approximate selectors, or a decoder
that directly combines several cells without exposing their line selectors.
All matrix weights below may depend on the individual selector.

\paragraph{Rank theorem used.}
Theorem 1.3 of Dvir, Saraf, and Wigderson,
\emph{Improved rank bounds for design matrices and a new proof of Kelly's theorem},
Forum of Mathematics, Sigma 2 (2014), e4,
\href{https://doi.org/10.1017/fms.2014.2}{doi:10.1017/fms.2014.2}, states that an
$m\times v$ complex matrix whose rows have at most $r$ nonzero entries,
columns have at least $k>0$ nonzero entries, and pairs of distinct columns
have at most $t$ common nonzero rows has
\begin{equation}\label{eq:dsw}
 \rank M\geq \frac{v}{1+r(r-1)mt/(vk^2)}.
\end{equation}
The theorem depends only on the zero pattern, so positive denominator products
need not coincide between selectors.

\begin{proposition}[Arbitrary retained supports]\label{prop:arbitrary}
Let $P\subseteq\FF_q^2$ have $v>0$ points, and let $\mathcal L$ be a family of
$m$ distinct nonvertical affine lines. Suppose every point of $P$ lies on at
least $k>0$ members of $\mathcal L$. Assign each $p\in P$ an address
$x_p\in\{0,1\}^n$. For each $\ell\in\mathcal L$, suppose
\[
 S_\ell(x)=\beta_\ell+\sum_{h=1}^{c_\ell}
 \frac{A_{\ell h}(x)}{B_{\ell h}(x)},\qquad c_\ell\leq c,
\]
where $\beta_\ell$ is constant, $A_{\ell h},B_{\ell h}$ are real affine
polynomials, and every $B_{\ell h}$ is nonzero on these addresses. Suppose
$S_\ell(x_p)\neq0$ if and only if $p\in\ell$. Then
\[
 D(n,c):=\sum_{j=0}^{\min(c,n)}\binom nj
 \ \geq\ \frac{v}{1+q(q-1)m/(vk^2)}.
\]
In particular, original-model denominators positive on the entire Boolean
cube satisfy the denominator hypothesis.
\end{proposition}
\begin{proof}
Multiply each selector by its own denominator product. The resulting
polynomial $P_\ell$ has degree at most $c$ and exactly the same nonzero
pattern on $P$. Multilinear reduction preserves its values on Boolean
addresses and does not increase degree. Thus the rows of the matrix
$M_{\ell p}=P_\ell(x_p)$ lie in an evaluation space of dimension at most
$D(n,c)$. Each line contains at most $q$ points, and two distinct points
lie on at most one nonvertical line. Apply \eqref{eq:dsw} with $r=q,t=1$.
No common denominator across rows was used.
\end{proof}

\begin{corollary}[Active slope classes]\label{cor:active}
If $P=\FF_q^2$ and $\mathcal L$ consists of all lines in $k$ slope classes,
where $1\leq k\leq q$, then
\[
 D(n,c)\geq\frac{q^2k}{k+q-1}\geq\frac{kq}{2}.
\]
Zero residue directions must be deleted before applying this result.
For arbitrary retained supports, Proposition~\ref{prop:arbitrary} applies
only after establishing a positive minimum incidence degree.
\end{corollary}
\begin{proof}
There are $m=kq$ rows, $v=q^2$ columns, and every column has exactly $k$
nonzero entries. Substitute these parameters in Proposition~\ref{prop:arbitrary}.
\end{proof}

\begin{proposition}[All decoder bits remove the zero-residue loophole]
Let $q=2^a$, $a\geq1$, and let $t_1<\cdots<t_q$ be scalar tokens in an
interval $I$. For each bit $b\in\{1,\ldots,a\}$, let
\[
 R_b(t)=\gamma_b+\sum_{\alpha\in S_b}
       \frac{\eta_{b\alpha}}{d_\alpha-t},
 \qquad 0\neq\eta_{b\alpha}\in\RR,
\]
with distinct poles outside $I$. Suppose the strict sign vectors of
$(R_b(t_j))_{b=1}^a$ are pairwise distinct. Write $k_b=|S_b|\leq q$.
For every $b$ and every slope in $S_b$, suppose exact selectors as in
Proposition~\ref{prop:arbitrary}, of at most $c$ heads each, are available
for all lines of that slope. Addresses may depend on $b$. Then
\begin{equation}\label{eq:budget}
 a D(n,c)\geq\frac{q(q-1)}2.
\end{equation}
\end{proposition}
\begin{proof}
Clearing $R_b$ gives a nonzero polynomial of degree at most $k_b$; its
denominator has constant sign on $I$. Each of the $q-1$ consecutive token
gaps contains a sign change in at least one bit, hence a root of the
corresponding polynomial. Roots in distinct gaps are distinct, and
therefore $q-1\leq\sum_b k_b$. For $k_b>0$, Corollary~\ref{cor:active}
gives $D(n,c)\geq k_bq/2$. For $k_b=0$ the same weak inequality is trivial.
Sum these inequalities over $b$.
\end{proof}

\paragraph{Accounting consequence, and its limit.}
For $1\leq c\leq n$, $D(n,c)\leq(en/c)^c$ (the bound is immediate for
$c>n/2$ from $D\leq2^n$). Since $q-1\geq q/2$, \eqref{eq:budget} implies
\[
 2\log q\leq c\log(en/c)+\log(4a).
\]
As $q=2^a$, this forces $c=\Omega(a/\log n)$ in the growing-$a$ regime.
Consequently, a compiler that assigns a uniform budget of $c$ heads to
\emph{each} of its $q^2$ line selectors cannot improve the gain
$a/c$ beyond $O(\log n)$ by increasing the field size. This last accounting
statement assumes the uniform $cq^2$ cost; the bound on a maximum selector
budget alone is not a lower bound on the sum of unequal selector costs.
It also grants multiplication of a selector by its memory feature for free;
such multiplication can violate affine numerator admissibility and needs
its own construction. Sharing heads across many selectors, nonuniform
budgets, directly decoded joint memories, and controlled approximation
remain outside this conclusion.

\paragraph{Formal pieces already available.}
The scalar sign-change argument is formalized in
\texttt{JointFamilyScalarCapacity.lean}. The exact fixed-feature incidence
rank obstruction is formalized in \texttt{JointFamilyIncidenceObstruction.lean}.
The external weighted-design theorem and its composition here have not been
formalized. The positive local compiler improvement in
\texttt{JointCellRankRepair.lean} is independent of this written obstruction.
\end{document}
